\section{Introduction}

Unified understanding and generation models that integrate both visual comprehension and image synthesis within a single architecture represent a significant step toward general-purpose multimodal intelligence. Unlike specialist models that perform either visual question answering or text-to-image synthesis in isolation, unified architectures share representations across both modalities, creating the potential for cross-modal synergy: visual comprehension can inform more faithful image synthesis by grounding generation in fine-grained scene understanding, while generative experience can enrich perceptual reasoning by exposing the model to the causal structure of visual content. Several recent architectures have demonstrated that a single model can achieve competitive performance on both understanding and generation benchmarks simultaneously~\citep{chen2025januspro, ge2024seedx, wang2024emu3}, establishing unified understanding and generation models as a viable paradigm. However, despite this architectural convergence, the initial pretraining stage optimizes a broad distributional objective over large-scale paired data, and the resulting models exhibit substantial headroom for post-training refinement on both capabilities.

A central challenge in closing this gap is the supervision bottleneck. The prevailing approach to post-training relies on curated, task-specific human annotations: visual question-answer pairs for understanding, aesthetic preference labels for generation, or both. This supervision-dependent paradigm faces three structural limitations that become increasingly restrictive as unified models scale. First, annotation cost scales linearly with data volume and must be replicated separately for each modality, creating a compounding bottleneck. Second, human-curated labels introduce systematic biases toward the annotator population's perceptual and linguistic priors, which can narrow the model's competence to well-represented data distributions and limit its ability to improve on out-of-distribution content. Third, and perhaps most critically for unified models, supervised signals for understanding and generation are typically designed in isolation, treating the two modalities as independent objectives rather than exploiting the mutual dependence between them. A post-training method that derives its learning signal entirely from the model's own behavior, rather than external annotations, would eliminate these bottlenecks and enable the model to improve continuously on any source of unlabeled images.

The key observation motivating our approach is that a unified model's own understanding and generation capabilities, though imperfect, provide sufficient signal for self-improvement when properly structured. Consider a visual question posed to the solver under multiple semantically equivalent but syntactically distinct prompt framings. If the solver consistently produces the same answer regardless of how the question is phrased, this robust agreement reflects genuine comprehension of the visual content rather than superficial pattern matching. If the answers instead disagree across framings, the disagreement precisely localizes the boundary of the model's current competence: the question probes knowledge that the model has partially acquired but not yet consolidated. This prompt-perturbed self-consistency signal requires no external labels, adapts naturally to the model's evolving capability level, and provides a principled foundation for autonomous self-improvement. Moreover, the solver's understanding capability provides a natural evaluation mechanism for generation quality: if the solver can accurately answer diagnostic questions about a real image but fails to answer the same questions about the model's own synthesized version of that scene, the generation has demonstrably failed to preserve the relevant visual content.

Recent methods have begun to address the supervision challenge through various forms of self-improvement and reinforcement-based optimization, but each introduces dependencies that limit its generality. Role-based self-play approaches decompose the model into interacting components trained through adversarial or cooperative dynamics~\citep{su2025unigame, han2026unicorn, tang2026endogenous}, but require auxiliary components such as dedicated judges, perturbers, or reward activators that introduce their own training requirements. Reconstruction-based methods exploit the internal gap between understanding and generation quality~\citep{yan2025understandinggeneration, han2025internalgap}, but assume specific architectural couplings such as shared encoder-decoder pathways or access to intermediate hidden representations. Reinforcement learning formulations cast post-training as policy optimization with task-specific or self-derived rewards~\citep{mao2025unirl, jiang2025coreinforcement, hong2025dualselfrewards}, but typically depend on supervised fine-tuning as an initialization or interleaved stage. A recurring pattern emerges: existing methods typically rely on at least one of (i)~curated supervised fine-tuning data, (ii)~external reward models or judge components, (iii)~architecture-specific mechanisms, or (iv)~single-scalar reward signals that capture only one dimension of quality. No existing method simultaneously eliminates all four dependencies while demonstrating portability across architecturally distinct model families.

In this work, we propose a self-evolving training framework that enables a unified model to autonomously improve both understanding and generation using only unlabeled images, without human annotations, curated instruction data, or external reward models at any stage. The framework decomposes a frozen pretrained backbone into three co-evolving roles through lightweight LoRA adapters~\citep{hu2022lora}: a \emph{proposer} that generates visual questions about unlabeled images, a \emph{solver} that answers them through the understanding pathway, and a \emph{generator} that synthesizes images through the generation pathway. The learning signal emerges from the self-consistency mechanism described above: the proposer learns to target the boundary of the solver's competence, generating progressively harder questions that keep the solver in its most productive learning regime. As the solver improves on harder questions, its evaluations of generated images become more discriminative, which in turn provides a stronger training signal for the generator. This creates a virtuous cycle in which understanding improvement naturally drives generation improvement through the solver's dual role as both learner and internal evaluator, without requiring any direct coupling between the generation loop and the proposer. % An overview of the framework is illustrated in Figure~\ref{fig:framework}.

Two key technical challenges arise in realizing this framework. The first concerns the self-consistency signal itself. When all solver responses agree, the sample-level entropy is identically zero regardless of whether the question was trivially easy or narrowly succeeded with substantial internal hesitation. Yet questions in this ambiguous regime, where the model arrives at the correct answer under all framings but does so with considerable uncertainty in its token-level predictions, are precisely the most valuable for curriculum learning. They sit at the frontier of the solver's competence, where slightly harder variants would cause disagreement. Without a mechanism to detect this internal uncertainty, the proposer cannot distinguish genuinely easy questions from those that probe the decision boundary, and therefore cannot maintain productive difficulty pressure. We address this with Solver Token Entropy (STE), a continuous difficulty measure derived from the token-level prediction uncertainty of a single greedy forward pass. STE extracts the maximum entropy over the solver's softmax distributions during answer generation, capturing the point of greatest internal hesitation. Combined with prompt-perturbed self-consistency, STE yields a two-component training signal that covers the full difficulty spectrum from trivial to unsolvable, enabling automatic curriculum learning without manual threshold design.

The second challenge concerns generation assessment. Evaluating the quality of synthesized images purely from the model's own understanding is inherently difficult because no single self-derived measure captures all dimensions of image quality. A generated image could answer all diagnostic questions about specific attributes correctly while still exhibiting holistic failures such as incoherent scene composition, context omission, or visual artifacts in regions that no question happened to probe. Conversely, a generated image could produce a caption whose overall meaning matches the original prompt while containing fine-grained factual errors in object color, count, or spatial arrangement that caption-level similarity overlooks. We address this by designing a multi-scale generation assessment mechanism that combines two complementary scoring functions. QA fidelity scoring probes specific semantic attributes by comparing the solver's answers about the real and generated images, catching fine-grained content errors. Cycle-consistent captioning evaluates holistic coherence by measuring whether the generated image elicits a description whose overall meaning aligns with the original generation prompt. This multi-component reward provides denser optimization signal than single-scalar alternatives and reduces the risk of reward gaming, as the generator must simultaneously satisfy both fine-grained factual consistency and global semantic alignment.

The entire framework operates at the level of role decomposition and cyclic self-supervised optimization, without dependence on architecture-specific components such as intermediate hidden representations, encoder-decoder symmetry, or reversible likelihood interfaces. We validate this model-agnostic design by implementing the same algorithm with identical hyperparameters on three unified models that employ fundamentally different generation paradigms: diffusion-based denoising, flow-matching with a variational decoder, and autoregressive discrete-token prediction. To the best of our knowledge, this is the first self-supervised post-training framework demonstrated across three architecturally distinct families of unified understanding and generation models.

\noindent Our contributions are as follows:
\begin{itemize}
    \item We propose a self-evolving training framework that improves both understanding and generation capabilities of unified models using only unlabeled images, eliminating the need for human annotations, curated instruction data, and external reward models.
    \item We introduce Solver Token Entropy (STE), a continuous difficulty measure extracted from token-level prediction uncertainty in a single greedy forward pass, which resolves the degenerate-signal limitation of self-consistency in the unanimous-agreement regime and enables effective curriculum learning from a cold start.
    \item We design a multi-scale generation assessment mechanism that combines fine-grained QA fidelity scoring with holistic cycle-consistent captioning, using the solver as an internal evaluator to couple understanding quality directly to the generation training signal.
    \item We demonstrate model-agnostic generality by applying the same algorithm with identical hyperparameters across three architecturally distinct unified models spanning diffusion-based, flow-matching, and autoregressive generation paradigms.
\end{itemize}

\noindent The remainder of this paper is organized as follows. Section~\ref{sec:related_work} surveys existing post-training methods and positions our approach within the landscape. Section~\ref{sec:method} presents the framework in detail. Section~\ref{sec:experiments} provides experimental validation across all three architectures, and Section~\ref{sec:conclusion} concludes with a discussion of limitations and future directions.
